\documentclass[10pt]{article}
\usepackage[T1]{fontenc}
\usepackage[utf8]{inputenc}
\usepackage[french]{babel}
\usepackage{amsmath,amssymb,mathrsfs,mathpazo}
\usepackage{geometry}
\geometry{a4paper,margin=2cm}
\usepackage{graphicx,array,tabularx,tabularray,colortbl,multirow,multicol,booktabs,enumitem,pifont,ulem,fancybox,diagbox,cancel}
\usepackage{pstricks,pstricks-add,pst-plot,pst-tree,pst-node,pst-text,pst-func,pst-math,pst-bspline,pst-3dplot,pst-eucl,pst-solides3d}
\let\pscalculate\undefined
\usepackage{pst-bezier}
\usepackage{tikz,pgfplots,tkz-tab}
\usetikzlibrary{arrows,arrows.meta,patterns,calc,positioning,trees,decorations.markings,shapes.geometric,shapes.misc,shapes.arrows,intersections,decorations.pathmorphing,decorations.pathreplacing,backgrounds,shadows,babel}
\pgfplotsset{compat=1.18}
\usepgfplotslibrary{fillbetween}
\usepackage[np]{numprint}
\usepackage{listings,makecell,pas-tableur,twoopt}
\let\barre\undefined
\usepackage[pstricks]{tablvar}
\definecolor{darkgreen}{RGB}{0,100,0}
\providecommand{\red}{\color{red}}
\providecommand{\blue}{\color{blue}}
\providecommand{\green}{\color{green}}
\providecommand{\black}{\color{black}}
\providecommand{\white}{\color{white}}
\providecommand{\gray}{\color{gray}}
\providecommand{\cyan}{\color{cyan}}
\providecommand{\magenta}{\color{magenta}}
\providecommand{\yellow}{\color{yellow}}
\providecommand{\violet}{\color{violet}}
\providecommand{\orange}{\color{orange}}
\providecommand{\euro}{\texteuro}
\providecommand{\EUR}{\texteuro}
\providecommand{\ieme}{\textsuperscript{e}}
\providecommand{\degres}{\ensuremath{^\circ}}
\providecommand{\degree}{\ensuremath{^\circ}}
\providecommand{\vv}[1]{\overrightarrow{#1}}
\providecommand{\Coord}[1]{#1}
\providecommand{\decosix}{}
\providecommand{\starredbullet}{\ensuremath{\bullet}}
\pagestyle{empty}
\setlength{\parindent}{0pt}
\usepackage{adjustbox}\DeclareUnicodeCharacter{2236}{:}\DeclareUnicodeCharacter{2190}{\ensuremath{\leftarrow}}
\providecommand{\R}{\mathbb{R}}
\providecommand{\N}{\mathbb{N}}
\providecommand{\Z}{\mathbb{Z}}
\providecommand{\Q}{\mathbb{Q}}
\providecommand{\D}{\mathbb{D}}
\providecommand{\C}{\mathbb{C}}
\providecommand{\e}{{\mathrm{e}}}
\providecommand{\vect}[1]{\overrightarrow{#1}}
\providecommand{\vectt}[1]{\overrightarrow{\mathrm{#1}}}
\providecommand{\vvt}[1]{\overrightarrow{\mathrm{#1}}}
\providecommand{\np}[2][]{{#2~#1}}
\providecommand{\vv}[1]{{\overrightarrow{#1}}}
\renewcommand{\d}{{\mathrm{d}}}
\providecommand{\degres}{{^\circ}}
\providecommand{\degre}{{^\circ}}
\providecommand{\blue}{\color{blue}}
\providecommand{\red}{\color{red}}
\providecommand{\enskip}{\quad}
\providecommand{\llbracket}{\lbrack\!\lbrack}
\providecommand{\rrbracket}{\rbrack\!\rbrack}
\providecommand{\euro}{{\text{€}}}
\providecommand{\no}{{\text{n°}}}
\providecommand{\interval}[2]{{[#1;#2]}}
\providecommand{\intervFO}[2]{{[#1;#2[}}
\providecommand{\dotfill}{\ldots}

\providecommand{\ds}{\displaystyle}
\providecommand{\textding}[1]{#1}
\providecommand{\Oij}{\(\mathrm{O};\vec{\imath},\vec{\jmath}\)}
\providecommand{\Oijk}{\(\mathrm{O};\vec{\imath},\vec{\jmath},\vec{k}\)}
\providecommand{\Ouv}{\(\mathrm{O};\vec{u},\vec{v}\)}
\providecommand{\intervFF}[2]{\ensuremath{[#1;#2]}}
\renewcommand{\euro}{\ifmmode\text{\texteuro}\else\texteuro\fi}
\DeclareUnicodeCharacter{00B0}{\ifmmode^{\circ}\else\textdegree\fi}
\DeclareUnicodeCharacter{20AC}{\ifmmode\text{\texteuro}\else\texteuro\fi}
% Énoncés : Éducation nationale. Transcription et corrigés : APMEP, Denis Vergès et contributeurs. Découpage et classement : CoopMaths, Sébastien Lozano, Éric Elter et contributeurs. Adaptation de la présentation : Maths Explorer.
% CC BY-SA : https://coopmaths.fr/www/about/
% https://coopmaths.fr/alea/static/e3c/2020/tex/e3c_2020_00_sujet51_3.tex
\begin{document}
\noindent\textit{APMEP / CoopMaths. CC BY-SA. Présentation adaptée par Maths Explorer.}\par\medskip
\medskip

Un journal hebdomadaire est sur le point d'être créé.

Une étude de marché aboutit à deux estimations différentes concernant le nombre de journaux vendus :
\begin{itemize}
\item  1\up{re} estimation : \np{1000} journaux vendus lors du lancement, puis une progression des ventes de 3\,\% chaque semaine.
\item 2\up{e} estimation : \np{1000} journaux vendus lors du lancement, puis une progression régulière de 40 journaux supplémentaires vendus chaque semaine.
 \end{itemize}
On considère les suites $\left(u_n\right)$ et $\left(v_n\right)$ telles que, pour tout entier naturel $n\geqslant 1$, $u_n$ représente le nombre de journaux vendus la $n$-ième semaine selon la première estimation et $v_n$ représente le nombre de journaux vendus la $n$-ième semaine selon la deuxième estimation. 

Ainsi, $u_1 = v_1 = \np{1000}$.

\medskip

\begin{enumerate}
\item On considère la feuille de calcul ci-dessous :

\begin{center}
\begin{tabularx}{0.5\linewidth}{|>{\columncolor{blue}}m{0.35cm}|*{3}{>{\centering \arraybackslash}X|}}
\hline
\rowcolor{blue}&A&B&\cellcolor{yellow}C\\\hline
1&$n$&$u_n$&$v_n$\\\hline
2&1&\np{1000}&\np{1000}\\\hline
3&2&\np{1030}&\np{1040}\\\hline
4&3&\np{1060.9}&\np{1080}\\\hline
5&4&\np{1092.727}&\np{1120}\\\hline
\end{tabularx}
\end{center}

Quelle formule, saisie en B3 et recopiée vers le bas, permet d'obtenir les termes de la suite $\left(u_n\right)$ ?

\medskip

\item 
	\begin{enumerate}
		\item  Donner la nature de la suite $(u_n)$ puis celle de la suite $(v_n)$. Justifier.
		\item Montrer que pour tout entier naturel $n\geqslant1$, $v_n=960+40n$.
		\item Écrire, pour tout entier naturel $n\geqslant1$, l'expression de $u_n$ en fonction de $n$.
	\end{enumerate}
\item On définit, pour tout entier $n\geqslant 1$, la suite $\left(w_n\right)$ par $w_n=v_n-u_n$.

On donne ci-dessous un extrait de son tableau de valeurs :

\begin{center}
\begin{tabular}[]{|*{3}{c|}m{1.2cm}|*{4}{c|}}\hline
$n$		&1&2	&	&19&20&21		&22\\\hline
$w_n$	&0&10	&	&18&6&$-6$	&$-20$\\\hline
\end{tabular}
\end{center}

À partir de quelle semaine le nombre de journaux vendus d'après la première estimation devient-il supérieur au nombre de journaux vendus d'après la deuxième estimation ?
\end{enumerate}

\end{document}
